\documentclass{article}
\usepackage{amsmath,amssymb}
\usepackage[margin=1in]{geometry}
\usepackage{url}
\title{The regional rank bound in ALVV}
\author{TNLean contributors}
\date{October 29, 2026}
\begin{document}
\maketitle

\paragraph{Source and disposition.}
Arad, Landau, Vazirani and Vidick, arXiv:1602.08828v2,
Definition 5.10 and Theorem 5.14(2), assert that the regional
frustration-free replacement has operator Schmidt rank at most $d^2+2$
at every cut. For an interior region, this is false even for eight qubits,
nearest-neighbor orthogonal projections, a unique frustration-free ground
state, and a positive spectral gap. The bound is therefore not formalized.
Issue \#9849 must remain open for the source correction and its disposition.
The independent kernel and gap statements are unaffected.

\paragraph{Local projections and the ground state.}
Number eight qubits $0,\ldots,7$ and put
\begin{align}
    z & =|00\rangle\langle00|,             &
    v & =|01\rangle+|10\rangle+|11\rangle, &
    Q & =z+\tfrac13|v\rangle\langle v|,    &
    P & =I-Q,                              & P_0 & =I-z.
    \label{eq:alvv_ff_example_local}
\end{align}
Let $h_0=(P_0)_{01}$ and $h_i=P_{i,i+1}$ for $1\le i\le6$.
All terms are orthogonal projections and annihilate $|0\rangle^{\otimes8}$.
The kernel is one-dimensional: $h_0\psi=0$ pins qubits $0,1$ to zero.
If qubit $i$ is already pinned, then $h_i\psi=0$ requires its two-site
coefficient to lie in $\operatorname{span}\{|00\rangle,v\}$.
The coefficients of $|10\rangle,|11\rangle$ are zero, so the coefficient
of $v$ vanishes; qubit $i+1$ is pinned as well. Induction pins every qubit.
This argument applies coefficientwise in every configuration of the other
qubits, so it includes arbitrary entangled vectors.
The positive operator $H=\sum_i h_i$ consequently has a unique zero-energy
state and a strictly positive gap in finite dimension. A lower gap
parameter can be chosen in $(0,1]$ if that normalization is required.

\paragraph{The actual regional replacement.}
Replace the bonds $1,\ldots,5$, corresponding to qubit region
$J=\{1,\ldots,6\}$. Both exterior bonds are retained. In lexicographic
qubit order, Definition 5.10 gives exactly
\begin{align}
    \widetilde H
     & =P_0\otimes I_{64}+I_{64}\otimes P+2I_{256}
    -I_4\otimes Q\otimes Q\otimes I_4
    -I_2\otimes Q\otimes Q\otimes Q\otimes I_2.
    \label{eq:alvv_ff_example_replacement}
\end{align}
At the cut $0123\mid4567$, reshuffle the entries into the coefficient
matrix $R_{16a+a',16b+b'}=\widetilde H_{16a+b,16a'+b'}$,
where $a,a',b,b'\in\{0,\ldots,15\}$.
Operator Schmidt rank equals the rank of $R$ over $\mathbb C$.

Select rows $(85,0,17,1,16,34,68)$.
Select columns $(170,0,8,128,34,136,153)$ in that order.
The resulting minor is $N/27$, with
\begin{align}
    N          & =\begin{pmatrix}
                      95 & 69 & -3 & -3 & 89 & 75 & 93 \\
                      60 & 0  & 0  & 0  & 36 & 36 & 54 \\
                      66 & 36 & -9 & -9 & 60 & 42 & 60 \\
                      -3 & 0  & 0  & -9 & 0  & -9 & -9 \\
                      -3 & 0  & -9 & 0  & 0  & -9 & -9 \\
                      68 & 36 & 0  & 0  & 60 & 48 & 66 \\
                      89 & 45 & 0  & 0  & 69 & 69 & 87
                  \end{pmatrix},     \\
    \det N     & =-30233088,                         &
    \det(N/27) & =-\frac{512}{177147}\ne0.
    \label{eq:alvv_ff_example_minor}
\end{align}
Thus $\operatorname{OSR}(\widetilde H)\ge7>6=d^2+2$.
In fact equality holds: the two exterior sums together with $2I$ have
rank at most two, the even regional product has rank one, and the odd
regional product has rank at most four. These are actual tensor
factorizations across the indicated cut.

\paragraph{Verification and missing source step.}
The standard-library-only script
\path{scripts/check_alvv_ff_rank_counterexample.py} computes the
49 selected entries directly from (\ref{eq:alvv_ff_example_replacement}),
checks every local symmetry and idempotence equation over $\mathbb Q$,
and evaluates both displayed determinants by the Leibniz formula over
exact fractions. No floating-point rank tolerance is used. This is
an exact arithmetic certificate accompanying the mathematical argument;
a full Lean counterexample theorem is not asserted.
The source proof counts the exterior contribution as rank one when a cut
lies inside $J$, although there can be nontrivial exterior terms on both
sides. It also counts a regional complement as though it were its product.
The elementary factorization gives $d^2+3$ for interior regions, whereas
$d^2+2$ remains consistent with one-sided or whole-chain regions.
The corrected all-cut bound $d^2+3$ and the prefix/suffix bound $d^2+2$
are now proved for the actual regional matrix replacement. These corrected
statements are not a formalization of the unrestricted printed bound.

\paragraph{Endpoint convention.}
The source numbers qudits from one and explicitly defines
$J=\{j_0,\ldots,j_1\}$ and
$H_J=\sum_{i=j_0}^{j_1-1}h_i$, with $h_i$ acting on qudits $i,i+1$.
Its subsequent description of $J_e,J_o$ as parity subsets of $J$ would
include $j_1$, which is a qudit index but not an index of a term in this
explicit sum. We use the parity subsets of the bond set
$\{j_0,\ldots,j_1-1\}$ throughout. Including $h_{j_1}$ in the
regional product while retaining it in the exterior sum would count that
boundary bond twice and would be a different construction.
In zero-based coordinates $a=j_0-1$, $b=j_1-1$, the site region is
$\lbrack a,b+1)$ and the selected bond interval is $\lbrack a,b)$.
The original even and odd color names swap under this index shift;
their sum and all rank bounds are unchanged. In the example the site
region is $\lbrack1,7)$ and precisely the bonds $1,\ldots,5$ are replaced.
For $m$ original bonds, a prefix has $a=0$ and a suffix has $b\ge m$.

\paragraph{Squared norms in the gap comparison.}
Theorem 5.14's proof must use squared norms on both sides.
For commutation degree $g$, detectability and the original gap give
$(\gamma+g^2)\|D\psi\|^2\le g^2$ on unit vectors orthogonal to the
full original kernel. The imported converse with constant three applies
to every unit vector, giving
$1-3\langle\psi,\widetilde H\psi\rangle\le\|D\psi\|^2$.
Consequently the quadratic-form gap is at least
$\gamma/[3(\gamma+g^2)]$. For $g=2$ and $0<\gamma\le1$, it is at least
$\gamma/15$. An eigenvector assumption is unnecessary with this converse.
The full comparison is now assembled for the actual finite chain: disjoint-bond
commutation, degree two after parity ordering, preservation of the complete
kernel, and equality with the actual regional replacement are all derived
from the original local terms. Together with the corrected rank estimates,
this proves the corrected version of the chain theorem without supplying
any of these intermediate conclusions as hypotheses.
\end{document}
